This work aims to help model providers defend against unauthorised distillation, and it also describes ways to defeat the defences it studies: padding colluder accounts with natural text, ramping a campaign through the calibration weeks, choosing which accounts the detector flags, and re-weighting real user prompts toward a useful region. We report them because a defender cannot evaluate a monitor without them, and because each follows directly from the monitor's design or from long-known ideas (mimicry and boiling-frog poisoning); documented attackers already go further by relaying real users' requests at volumes far beyond ours. We will release the detectors, the calibration protocol and the census script openly, and the campaign constructions (rewriting prompts, scheduling and ramp code) to researchers and providers on request.

\emph{Data.} WildChat-1M and WildChat-4.8M are public under the AI2 ImpACT license, with accounts released as hashed IP addresses. We did not attempt to re-identify any user, inspected conversations only through automated processing, and reproduce no user prompt; every result is an aggregate over windows of tens to thousands of accounts.

\emph{Named organisations.} We name companies only as providers' and governments' public reports name them, describe their conduct as reported allegations, and make no independent claim about them.

\emph{Harms of deployment.} Population-level monitoring pools all users' queries into one statistic and can open side channels \citep{debenedetti2024privacy}; the monitor flags legitimate evaluation pipelines as readily as distillation; and randomised throttles degrade service for any legitimate account in a flagged group. Enforcement should follow human review, under the governance any population-level telemetry needs.
